\documentclass[11pt,a4paper]{article}
\usepackage[T1]{fontenc}
\usepackage[utf8]{inputenc}
\usepackage{amsmath,amssymb,amsthm}
\usepackage[margin=1in]{geometry}
\usepackage[colorlinks=true,linkcolor=blue,urlcolor=blue]{hyperref}
\theoremstyle{plain}
\newtheorem{theorem}{Theorem}
\newtheorem{lemma}[theorem]{Lemma}
\newtheorem{proposition}[theorem]{Proposition}
\newtheorem{corollary}[theorem]{Corollary}
\theoremstyle{definition}
\newtheorem{remark}[theorem]{Remark}

\title{The empirical recurrence for OEIS A184692, proved by transfer matrix}
\author{Adrian Perez Fontelles\\ \small Independent researcher}
\date{4 September 2026}
\begin{document}
\maketitle

\begin{abstract}
OEIS A184692 counts arrays of width $7$ over $\{0,1,2\}$ in which no cell equals both the cell above it and the cell to its left. The entry carries an
empirical recurrence of order $36$ contributed by R.~H.~Hardin. It is true, and it is
decidable rather than empirical: what a cell asks reaches one row above
it and none below, so the condition on a whole row is settled by
$2$ consecutive rows and the count is a walk count on $S=2188$ states.
\end{abstract}

\noindent\small 2020 Mathematics Subject Classification. 05A15, 05B20, 15A18.\normalsize

\section{The conjecture}

OEIS A184692 is ``1/3 the number of nX7 0..2 arrays with no element equal both to the element above and to the element to its left.''. It has offset $1$ and begins
\[
729, 786432, 831474432, 878052505824, 927173138866560, 979037502950237760, 1033802803683039013248, 1091631544523681779567680,\ \dots
\]
The entry states, as a conjecture contributed by its author and never marked settled:
\begin{quote}
Empirical: a(n)=1373*a(n-1)-382385*a(n-2)+54747930*a(n-3)-5050986386*a(n-4)+332267679041*a(n-5)-16472142082703*a(n-6)+636007583672478*a(n-7)-19508251954601454*a(n-8)+480569383283952657*a(n-9)-9546208348531626993*a(n-10)+152598952268814061851*a(n-11)-1946846244225175080346*a(n-12)+19518911243309328077698*a(n-13)-150002896733353526035184*a(n-14)+849776274453519925228352*a(n-15)-3327166677806732426576264*a(n-16)+7899544532604413646690344*a(n-17)-6508673827964553701633600*a(n-18)-20940104457639861466791296*a(n-19)+72463799516214011378881024*a(n-20)-61429017906552923870936064*a(n-21)-129536619155578798140518400*a(n-22)+380737956559758296779800576*a(n-23)-211491273596407625450717184*a(n-24)-649598231455057371522924544*a(n-25)+1411628961483792374550233088*a(n-26)-535856510128456063711707136*a(n-27)-2479363142583076923579039744*a(n-28)+5970652384046936507902066688*a(n-29)-7499573050858256420655071232*a(n-30)+6280307951272279750928760832*a(n-31)-3675985493399257320121171968*a(n-32)+1481320042982141475269115904*a(n-33)-383743024338594997911683072*a(n-34)+54171200367877184095780864*a(n-35)-2546711369956156422225920*a(n-36)
\end{quote}
As of the ``Last modified'' line on the live entry (Oct 03 2025 02:51:00, revision 11) this is still
recorded as empirical, and nothing on the entry records it as proved.

\section{What is being counted}

The arrays have $7$ columns and a growing number of rows, with entries in $\{0,1,2\}$. Fix the
list of offsets the entry names,
\[
\Delta=\{\,(-1,0), (0,-1)\,\},
\]
and for a cell $(i,j)$ write
\[
S(i,j)=\{\,(i+\delta_1,\,j+\delta_2) : (\delta_1,\delta_2)\in\Delta\,\}\cap
\text{(the cells of the array)}
\]
for the neighbours it actually has. Only offsets landing inside the array count, so a cell on
the border simply has fewer neighbours.

The condition is that no cell equals BOTH the cell above it and the cell to its left:
\[
\neg\bigl(x(i,j)=x(i-1,j)\ \wedge\ x(i,j)=x(i,j-1)\bigr).
\]
A cell in the first row or the first column is missing one of the two and so cannot equal both; it is under no condition. The entry publishes $1/3$ of that count; the division is by the constant 3 and changes nothing about the recurrence.

\section{A window of 2 rows}

The offsets reach one row up and none down, so everything the condition asks of one row is
decided by that row together with one row above it and none below --- a window of
$2$ consecutive rows, and nothing wider. Take the array one row at a time
and let the state be the last row, with a distinguished start standing for the
array before its first row, so that a walk of $n$ steps is an array of $n$ rows.

Appending a row completes exactly one window, and the row it settles is the bottom row of that window; a step is
allowed when that row passes. Each row of the finished
array is therefore tested exactly once.

Every state may end an array, so $\tau=\mathbf{1}$.
Thus
\[
a(n)\;=\;\iota^{\!\top}M^{\,j}\tau ,\qquad S=2188,
\]
for the transfer matrix $M$ on those $S$ states. The walk index $j$ and the entry's own $n$
differ by a constant, read off by matching the published terms rather than assumed. In
particular $a$ is $C$-finite.

\section{The proof}

\begin{lemma}\label{lem:crit}
Let $q(t)=t^{36}-\sum_i c_it^{36-i}$ be the characteristic polynomial of the
conjectured recurrence. The residual $a(n)-\sum_ic_ia(n-i)$ equals
$\iota^{\!\top}M^{\,j}q(M)\tau$ for the corresponding walk index $j$, so the recurrence
holds from some point on if and only if $\iota^{\!\top}M^{j}q(M)\tau=0$ for all large $j$,
and it is enough to examine $j<S$.
\end{lemma}

\begin{proof}
Factor $M^{j}$ out of $M^{j+36}-\sum_ic_iM^{j+36-i}$. For the bound, $M^{S}$
is an integer combination of $I,M,\dots,M^{S-1}$ by Cayley--Hamilton, so the numbers
$u_j=\iota^{\!\top}M^{j}q(M)\tau$ obey the monic recurrence given by the characteristic
polynomial of $M$; $S$ consecutive zeros force every later one, and the last nonzero $u_j$
pins the threshold exactly.
\end{proof}

\begin{theorem}
\[
a(n)\;=\;1373\,a(n-1) - 382385\,a(n-2) + 54747930\,a(n-3) - 5050986386\,a(n-4) + 332267679041\,a(n-5) - 16472142082703\,a(n-6) + 636007583672478\,a(n-7) - 19508251954601454\,a(n-8) + 480569383283952657\,a(n-9) - 9546208348531626993\,a(n-10) + 152598952268814061851\,a(n-11) - 1946846244225175080346\,a(n-12) + 19518911243309328077698\,a(n-13) - 150002896733353526035184\,a(n-14) + 849776274453519925228352\,a(n-15) - 3327166677806732426576264\,a(n-16) + 7899544532604413646690344\,a(n-17) - 6508673827964553701633600\,a(n-18) - 20940104457639861466791296\,a(n-19) + 72463799516214011378881024\,a(n-20) - 61429017906552923870936064\,a(n-21) - 129536619155578798140518400\,a(n-22) + 380737956559758296779800576\,a(n-23) - 211491273596407625450717184\,a(n-24) - 649598231455057371522924544\,a(n-25) + 1411628961483792374550233088\,a(n-26) - 535856510128456063711707136\,a(n-27) - 2479363142583076923579039744\,a(n-28) + 5970652384046936507902066688\,a(n-29) - 7499573050858256420655071232\,a(n-30) + 6280307951272279750928760832\,a(n-31) - 3675985493399257320121171968\,a(n-32) + 1481320042982141475269115904\,a(n-33) - 383743024338594997911683072\,a(n-34) + 54171200367877184095780864\,a(n-35) - 2546711369956156422225920\,a(n-36)
\]
for every $n>36$.
\end{theorem}

\begin{proof}
The vector $w=q(M)\tau$ was formed by $36$ matrix--vector products in exact integer
arithmetic and $\iota^{\!\top}M^{j}w$ evaluated for $j=0,1,\dots$, again exactly; those
integers vanish from the index corresponding to $n=36+1$ onwards and the run continued
until the working vector was identically zero, which settles every larger $j$ at once.
\end{proof}

\section{Verification}

Four checks, each independent of the algebra above.

First, the model reproduces all $10$ terms the entry publishes, exactly, in integer
arithmetic.

Second, the count was recomputed for the smallest arrays straight from the definition: every
array of width $7$ over $\{0,1,2\}$ with a few rows was written out cell by cell and each cell
tested against the neighbours it has. No window, no state and no end vector are involved in that
computation, and the published terms came back.

Third, the conjectured recurrence was evaluated directly on the published terms, with no
matrices involved, and holds wherever the proved range and the data overlap.

Fourth, the annihilation test of Lemma~\ref{lem:crit} was carried out in exact integer
arithmetic on the whole vector rather than on a sampled prefix.

\begin{thebibliography}{9}
\bibitem{oeis} The OEIS Foundation, \emph{The On-Line Encyclopedia of Integer
Sequences}, \url{https://oeis.org}, 2026. Sequence A184692.
\bibitem{stanley} R.~P.~Stanley, \emph{Enumerative Combinatorics}, Volume 1, 2nd ed.,
Cambridge University Press, 2011. (Transfer-matrix method, Section 4.7.)
\bibitem{horn} R.~A.~Horn and C.~R.~Johnson, \emph{Matrix Analysis}, 2nd ed., Cambridge
University Press, 2013.
\end{thebibliography}

\end{document}
